\section{Conductivité électronique d’un spinelle : la magnétite \ce{Fe3O4}}
%Marucco, ex12.3

\setcounter{equation}{0}
\setcounter{Q}{0}

La magnétite \ce{Fe3O4}, de structure spinelle, contient des ions \ce{Fe^{2+}} et \ce{Fe^{3+}}. La conductivité de ce matériau, à température ambiante, est relativement élevée ($200 ~ \si{\per\ohm\per\centi\meter}$). La mobilité des porteurs de charges est peu élevée ($\approx 1 ~ \si{\centi\meter\square\per\volt\per\centi\meter}$) et le libre parcours moyen de ces porteurs correspond à la distance \ce{Fe}-\ce{Fe}.

\Q Préciser l’état de spin des ions fer dans ce matériau. Le fer a pour structure électronique $3d^6 4s^2$.

\Q Proposer un modèle de conduction pour expliquer les propriétés électroniques de la magnétite. En effet, la faible valeur de la mobilité des porteurs et les propriétés magnétiques, liées à des états localisés de ces porteurs, sont en contradiction avec un modèle de bandes.
